\documentclass{siam/siamart}

\usepackage{amsmath, amssymb}
\usepackage{booktabs}
\usepackage{microtype}
\usepackage{xcolor}
\usepackage{listings}
\usepackage{graphicx}
\usepackage{algpseudocode}

\newsiamremark{remark}{Remark}
\crefname{remark}{Remark}{Remarks}
% Fix siamart.cls's refstepcounter override: it strips \@currentcounter,
% which TeX Live 2026's kernel and cleveref's First Aid patch both require.
\makeatletter
\def\cref@old@refstepcounter#1{%
  \stepcounter{#1}%
  \edef\reserved@a{#1}%
  \ifx\reserved@a\ltx@star@counter\else
    \let\@currentcounter\reserved@a
  \fi
  \protected@edef\@currentlabel{%
    \csname p@#1\expandafter\endcsname\csname the#1\endcsname}%
}%
\makeatother

% \lstdefinestyle{python}{
%   language=Python,
%   basicstyle=\ttfamily\small,
%   keywordstyle=\color{blue!70!black}\bfseries,
%   commentstyle=\color{gray},
%   stringstyle=\color{orange!80!black},
%   showstringspaces=false,
%   frame=single,
%   framesep=4pt,
%   xleftmargin=8pt,
%   xrightmargin=8pt,
%   belowskip=6pt,
%   aboveskip=6pt,
% }

\title{Krylov Methods for the Markowitz Portfolio}

\author{Thomas Schmelzer%
  \thanks{Jebel Quant Research (\email{thomas@jqr.ae}).}
  \and
  Raphael Hauser%
  \thanks{Mathematical Institute, University of Oxford.}
  \and
  Stephen Boyd%
  \thanks{Department of Electrical Engineering, Stanford University.}
}

\headers{Krylov Methods for the Markowitz Portfolio}{T.~Schmelzer, R.~Hauser, and S.~Boyd}

\begin{document}

\maketitle

\begin{abstract}
The Markowitz mean-variance portfolio problem, subject to general linear equality and
inequality constraints, admits a symmetric indefinite KKT saddle-point formulation that
MINRES \cite{paige1975} and a constraint-eliminated variant of CG \cite{hestenes1952}
approximate efficiently---without ever forming the sample covariance matrix.
Extending the minimum-variance Krylov framework of \cite{minvar_paper}, the
present paper identifies four generalisations required for the full Markowitz setting:
general linear equality constraints encoded in a matrix $A$ (factor-exposure targets,
sector-weight targets); general linear inequality constraints $C^\top w \leq d$ (e.g.\ sector caps)
handled by an active-set loop; a return tilt $\rho\,\mu^\top w$; and a Ledoit-Wolf
ridge $\alpha(\|X\|_F^2/n)\|w\|^2$.  In each case the KKT saddle-point structure is preserved and
both solvers apply without modification.  The saddle-point operator is implemented
matrix-free as a \texttt{LinearOperator}, avoiding the $O(Tn^2)$ formation cost and
$O(n^2)$ storage of $X^\top X$; the active-set loop is decoupled from the inner solver
via a function-pointer interface, so the same loop drives the direct KKT, MINRES, and
conjugate gradient \cite{hestenes1952} solvers uniformly.

Numerical experiments on synthetic data and on a real S\&P~500 universe demonstrate
order-of-magnitude speedups over open-source interior-point solvers for single-point
and efficient-frontier solves, with and without Ledoit-Wolf shrinkage.
\end{abstract}

\section{Introduction}

Mean-variance portfolio optimisation, introduced by Markowitz \cite{markowitz1952}, remains the
dominant framework for quantitative asset allocation.  In practice, portfolios are subject to a
range of linear constraints beyond the basic budget requirement: sector caps, factor-exposure
limits, individual position bounds, and long-only restrictions are routinely imposed.  Solving the resulting
constrained quadratic programme efficiently is therefore of direct practical importance, especially
as the number of assets $n$ and the frequency of re-optimisation grow.

The companion paper \cite{minvar_paper} showed that the long-only \emph{minimum variance}
portfolio reduces to a symmetric indefinite KKT saddle-point system, and that MINRES
\cite{paige1975} and a constraint-eliminated variant of CG \cite{hestenes1952} approximate
its solution efficiently, yielding large speedups over CVXPY \cite{diamond2016}
on standard benchmarks.
The present paper extends this framework to the full \emph{Markowitz mean-variance} problem.
Four generalisations are required: (i) the single budget constraint is replaced by a general
equality constraint matrix $A$; (ii) general linear inequalities $C^\top w \leq d$ are
handled by an active-set loop that promotes binding constraints to equalities; (iii) a linear
return tilt $\rho\,\mu^\top w$ is added to the objective; and (iv) a Ledoit-Wolf \cite{ledoit2004} ridge
$\alpha(\|X\|_F^2/n)\|w\|^2$, with $\alpha\in[0,1]$.  In each case the KKT
saddle-point structure is preserved and MINRES applies without modification.

Krylov subspace methods for portfolio optimisation have previously been proposed by
\cite{bajeux2012}, who project the sample covariance matrix onto a Krylov subspace and
solve the resulting low-dimensional problem, requiring $X^\top X$ to be formed explicitly;
the present approach instead exploits the KKT saddle-point structure of the first-order
optimality conditions, applying MINRES to the indefinite system without ever forming
$X^\top X$.

A central contribution of the paper is a modular software architecture that separates concerns
cleanly.  The full problem specification---returns matrix, constraints, shrinkage parameters, and
return-tilt weight---is encapsulated in a single frozen \texttt{Problem} dataclass (\S\ref{sec:problem}).
The saddle-point operator is implemented matrix-free as a
\texttt{scipy.sparse.linalg.LinearOperator} (\S\ref{sec:solvers}), avoiding the $O(Tn^2)$
formation cost and $O(n^2)$ storage of $X^\top X$.  The active-set loop is decoupled from the
inner linear solver via a function-pointer interface (\S\ref{sec:activeset}), so the same loop
drives the direct KKT solve, MINRES, and conjugate gradient \cite{hestenes1952} solvers
without code duplication.

The remainder of the paper is organised as follows.  Section~\ref{sec:problem} states the
problem and introduces the \texttt{Problem} dataclass.  Section~\ref{sec:kkt} derives the KKT
saddle-point system, establishes nonsingularity and optimality, and analyses MINRES convergence
via the spectral structure of the saddle-point matrix.  Section~\ref{sec:activeset} describes
the active-set method.  Section~\ref{sec:solvers} presents the four solvers and their
implementation, including a discussion of the fundamental differences from CVXPY and clarabel.
Section~\ref{sec:results} reports numerical experiments.

\section{Problem Formulation}\label{sec:problem}

Let $X \in \mathbb{R}^{T \times n}$ be the demeaned (column-wise) daily return matrix for
$n$ assets over $T$ trading days, $\mu \in \mathbb{R}^n$ a vector of expected returns,
$\rho \geq 0$ a return-tilt weight, and $\alpha \in [0,1]$ a Ledoit-Wolf shrinkage intensity.  The general mean-variance
problem is
\begin{equation}\label{eq:primal}
\min_{w} \; (1-\alpha)\|Xw\|^2 + \frac{\alpha\|X\|_F^2}{n}\|w\|^2 - \rho\,\mu^\top w
\quad \text{subject to} \quad
A^\top w = b, \quad C^\top w \leq d
\end{equation}
where $A \in \mathbb{R}^{n \times m}$, $b \in \mathbb{R}^m$ encode $m$ linear equality constraints
(at minimum the budget constraint, and optionally factor-exposure targets, sector-weight mandates, etc.),
and $C \in \mathbb{R}^{n \times p}$, $d \in \mathbb{R}^p$ encode $p$ linear inequality constraints
(sector caps and long-only being the typical cases; the long-only requirement $w \geq 0$ is the special case $C = -I$, $d = \mathbf{0}$).
The parameter $\alpha \in [0,1]$ in~\eqref{eq:primal} is user-supplied: $\alpha=0$
recovers the unregularised sample covariance while $\alpha=1$ replaces it with an
isotropic target.  Following Ledoit and Wolf~\cite{ledoit2004}, a natural estimator
shrinks the sample covariance $X^\top X / T$ toward a scaled identity,
\[
  (1-\alpha)\,\frac{X^\top X}{T} + \alpha\,\frac{\|X\|_F^2}{Tn}\,I.
\]
The Ledoit-Wolf (2004) oracle minimises the expected Frobenius-norm loss to the
true covariance and gives a data-adaptive $\alpha$~\cite{ledoit2004}; it does
\emph{not} in general equal a fixed ratio.  However, for the specific case where
the true covariance is proportional to the identity (i.i.d.\ Gaussian data, as in
the Marchenko--Pastur benchmark of \S\ref{sec:results}), the oracle shrinkage
intensity converges to $\alpha^* = n/(n+T)$ as $n, T \to \infty$ at fixed ratio
$n/T$.  We therefore adopt $\alpha = n/(n+T)$ as a recommended default that is
exact for the synthetic benchmark and serves as a computationally simple
approximation for general data; other choices are valid and the data-adaptive
formula of~\cite{ledoit2004} should be used when the covariance structure is
unknown.  Multiplying by $T$, the portfolio variance term becomes
$(1-\alpha)\|Xw\|^2 + (\alpha\|X\|_F^2/n)\|w\|^2$ as in~\eqref{eq:primal}; the
solver computes the ridge coefficient $\alpha\|X\|_F^2/n$ from $X$ and $\alpha$
internally, at $O(Tn)$ cost via $\|X\|_F^2 = \sum_{t,i}X_{ti}^2$.
Setting $\rho = 0$, $\alpha = 0$, $A = \mathbf{1}$, $b = 1$, $C = -I$, $d = \mathbf{0}$
recovers the minimum variance problem of \cite{minvar_paper}.

\begin{remark}
The parameter $\rho$ is the risk-aversion coefficient: larger $\rho$ tilts
the portfolio toward higher expected return at the cost of higher variance.  In
practice $\rho$ is either set by the investor's risk preference or swept over a range
$[0, \rho_{\max}]$ to trace the efficient frontier, as in \S\ref{sec:results}.
\end{remark}

\begin{remark}
When $T < n$ and $\alpha = 0$ the matrix $X^\top X$ is rank-deficient, the objective
is not strictly convex, and the KKT system is singular; all solvers fail.  Setting
$\alpha > 0$ via Ledoit-Wolf shrinkage is therefore essential in the undersampled
regime $T < n$, and recommended in general.
\end{remark}

The companion package \cite{fast_minimum_variance} encapsulates the full
problem specification in a single immutable object; Table~\ref{tab:params}
lists its parameters.  Omitting a constraint matrix activates the long-only
minimum-variance defaults.  Each solver returns $(w, k)$ where $k$ is the
cumulative inner iteration count.  The Lagrange multipliers $\lambda$ (dual
variables, encoding the shadow price of each equality constraint) are available
from the KKT solution vector but are not surfaced in the current API; they
carry direct economic meaning as the marginal cost of tightening a constraint.

\begin{table}[ht]
\centering
\begin{tabular}{lll}
\toprule
Parameter & Type & Description \\
\midrule
$X$       & $\mathbb{R}^{T \times n}$ & Demeaned return matrix \\
$\rho$    & $\mathbb{R}_{\geq 0}$     & Return-tilt weight (default $0$) \\
$\alpha$  & $[0,1]$                   & Ledoit-Wolf shrinkage intensity (default $0$) \\
$\mu$     & $\mathbb{R}^n$            & Expected returns (default $\mathbf{0}$) \\
$A$       & $\mathbb{R}^{n \times m}$ & Equality LHS (default $\mathbf{1}$) \\
$b$       & $\mathbb{R}^m$            & Equality RHS (default $1$) \\
$C$       & $\mathbb{R}^{n \times p}$ & Inequality LHS (default $-I$) \\
$d$       & $\mathbb{R}^p$            & Inequality RHS (default $\mathbf{0}$) \\
\bottomrule
\end{tabular}
\caption{Parameters of the \texttt{Problem} specification.  Defaults reproduce
the long-only minimum-variance problem of \cite{minvar_paper}.}
\label{tab:params}
\end{table}

\section{The KKT System}\label{sec:kkt}

\subsection{Equality-constrained subproblem}


Ignoring the inequality constraints, the Lagrangian is
\[
\mathcal{L}(w, \lambda) = (1-\alpha)\|Xw\|^2 + \tfrac{\alpha\|X\|_F^2}{n}\|w\|^2 - \rho\,\mu^\top w - \lambda^\top(A^\top w - b).
\]
The stationarity condition $\nabla_w \mathcal{L} = 0$ gives
$2\bigl((1-\alpha)X^\top X + \tfrac{\alpha\|X\|_F^2}{n}I\bigr)w - A\lambda = \rho\mu$.
Combined with $A^\top w = b$, this yields the $(n+m)\times(n+m)$ saddle-point system
\begin{equation}\label{eq:kkt}
\begin{pmatrix} 2\bigl((1-\alpha)X^\top X + \tfrac{\alpha\|X\|_F^2}{n}I\bigr) & A \\ A^\top & 0 \end{pmatrix}
\begin{pmatrix} w \\ \lambda \end{pmatrix}
=
\begin{pmatrix} \rho\mu \\ b \end{pmatrix}.
\end{equation}
The matrix is \textbf{symmetric but indefinite}; MINRES \cite{paige1975} handles this directly.
Setting $\rho = 0$, $\alpha = 0$ recovers the minimum variance KKT system of \cite{minvar_paper} with the
right-hand side $(\mathbf{0},\, b)$. The constraint-eliminated CG formulation generalises analogously.

\subsection{Nonsingularity and optimality}

\begin{proposition}\label{prop:nonsingular}
Suppose $A$ has full column rank and $(1-\alpha)X^\top X + (\alpha\|X\|_F^2/n)I \succ 0$.  Then the
saddle-point matrix in~\eqref{eq:kkt} is nonsingular, and any solution
$(w^*, \lambda^*)$ of system~\eqref{eq:kkt} satisfies the first-order KKT
optimality conditions for problem~\eqref{eq:primal}; these conditions are
sufficient for global optimality.  The same holds for each active-set
subproblem with $A$ replaced by $A_{\mathcal{A}}$ and $b$ by $b_{\mathcal{A}}$,
provided $A_{\mathcal{A}}$ has full column rank.
\end{proposition}

\begin{proof}
The objective $(1-\alpha)\|Xw\|^2 + (\alpha\|X\|_F^2/n)\|w\|^2 - \rho\,\mu^\top w$ is strictly
convex whenever $(1-\alpha)X^\top X + (\alpha\|X\|_F^2/n)I \succ 0$, which holds for all $\alpha > 0$
and, when $\alpha = 0$, whenever $X$ has full column rank.  The equality constraints
are affine and $A$ has full column rank by assumption; since the gradients of the
active equality constraints are precisely the columns of $A$, full column rank of $A$
is equivalent to LICQ.  Combined with strict convexity, the KKT conditions are both
necessary and sufficient for global optimality \cite[Ch.~5]{boyd2004}.  The same
argument applies to each active-set subproblem: full column rank of $A_{\mathcal{A}}$
is equivalent to LICQ for the augmented equality-constrained subproblem, ensuring
necessity and sufficiency of the KKT conditions there as well.

Nonsingularity of the saddle-point matrix follows from the standard theory of
symmetric saddle-point systems \cite{benzi2005}: the $(1,1)$ block
$Q = 2((1-\alpha)X^\top X + (\alpha\|X\|_F^2/n)I)$ is positive definite and $A$ has full column rank, so
the matrix is nonsingular with $n$ positive and $m$ negative eigenvalues.
\end{proof}

\begin{remark}
In practice $\alpha > 0$ is always recommended: Ledoit-Wolf shrinkage
simultaneously ensures $(1-\alpha)X^\top X + (\alpha\|X\|_F^2/n)I \succ 0$
and compresses the spectrum, reducing the MINRES iteration count substantially.
\end{remark}

\subsection{Spectral analysis and MINRES convergence}\label{sec:spectrum}

The convergence of MINRES on system~\eqref{eq:kkt} is governed by the eigenvalue
distribution of $K$.  Since $Q = 2((1-\alpha)X^\top X + (\alpha\|X\|_F^2/n)I) \succ 0$ and $A$ has full column
rank $m$, the matrix $K$ has exactly $n$ positive and $m$ negative eigenvalues
\cite{benzi2005}.  Specifically, the positive eigenvalues lie in
\[
  \bigl[\,\lambda_{\min}(Q),\;\;
         \lambda_{\max}(Q) + \sigma_{\max}^2(A)\,/\,\lambda_{\min}(Q)\,\bigr]
\]
and the negative eigenvalues in $[-\sigma_{\max}^2(A)/\lambda_{\min}(Q),\;
-\sigma_{\min}^2(A)/\lambda_{\max}(Q)]$ \cite{benzi2005}.  When
$\sigma_{\max}^2(A) \ll \lambda_{\min}(Q)\,\lambda_{\max}(Q)$---satisfied for the
budget constraint $A = \mathbf{1}$ whenever $\alpha > 0$---the effective spectral
condition number is
\begin{equation}\label{eq:kappa}
  \kappa_{\mathrm{eff}} \;\approx\;
  \frac{(1-\alpha)\sigma_1^2(X) + \alpha\|X\|_F^2/n}{(1-\alpha)\sigma_n^2(X) + \alpha\|X\|_F^2/n}
  \;=\; \frac{\lambda_{\max}(Q)}{\lambda_{\min}(Q)},
\end{equation}
where $\sigma_1(X) \geq \cdots \geq \sigma_n(X)$ are the singular values of $X$.
Note that $\kappa_{\mathrm{eff}}$ is the condition number of the positive-definite
$(1,1)$ block $Q = 2((1-\alpha)X^\top X + (\alpha\|X\|_F^2/n)I)$, not of the full indefinite matrix $K$.
From the polynomial perspective of \cite{fischer1996}, Krylov methods approximate
the action of a matrix inverse by a low-degree polynomial evaluated at the matrix;
for an SPD system convergence is governed by how well $1/x$ can be approximated on
the spectrum, and the Chebyshev bound \cite{golub2013,fischer1996}
\begin{equation}\label{eq:minres_conv}
\frac{\|r_k\|}{\|r_0\|} \;\lesssim\; 2\!\left(\frac{\sqrt{\kappa}-1}{\sqrt{\kappa}+1}\right)^{k}
\end{equation}
quantifies this in terms of the condition number $\kappa$.  Small eigenvalues make
$1/x$ hard to approximate and slow convergence; the identity term $(\alpha\|X\|_F^2/n)I$
lifts the spectrum away from zero, reduces $\kappa$, and thereby improves polynomial
approximability.  In this sense Ledoit-Wolf shrinkage acts as a spectral
preconditioner that directly accelerates CG and, by extension, MINRES.
Applying this to the SPD $(1,1)$ block $Q$ with $\kappa = \kappa_{\mathrm{eff}}$
from~\eqref{eq:kappa}, and noting that under the stated condition
$\sigma_{\max}^2(A) \ll \lambda_{\min}(Q)\lambda_{\max}(Q)$ the indefinite $K$ is
spectrally dominated by $Q$, gives the practical proxy
\begin{equation}\label{eq:minres_conv_proxy}
\frac{\|r_k\|}{\|r_0\|} \;\lesssim\; 2\!\left(\frac{\sqrt{\kappa_{\mathrm{eff}}}-1}{\sqrt{\kappa_{\mathrm{eff}}}+1}\right)^{\!\lfloor k/2\rfloor},
\end{equation}
where the factor $\lfloor k/2 \rfloor$ (versus $k$ for CG on an SPD system) reflects
the two-sided spectrum of $K$; see \cite{greenbaum1997,choi2006} for convergence
analysis of MINRES on symmetric indefinite systems.
As active sector-cap columns enter $A_{\mathcal{A}}$ during the active-set loop,
$\sigma_{\max}(A_{\mathcal{A}})$ grows and this condition may weaken, making
$\kappa_{\mathrm{eff}}$ a less accurate proxy for later active-set steps.

\begin{remark}\label{rem:marchenko}
Ledoit-Wolf shrinkage reduces $\kappa_{\mathrm{eff}}$ directly.  For $T = 2n$
i.i.d.\ Gaussian data the Marchenko--Pastur spectral edges are
$\lambda_\pm = T(1 \pm 1/\sqrt{2})^2$ (eigenvalues of $X^\top X$; the corresponding
edges of $X^\top X/T$ are $(1 \pm 1/\sqrt{2})^2$), giving
$\sigma_1^2(X) \approx 5.83n$, $\sigma_n^2(X) \approx 0.17n$, and
$\mathbb{E}\|X\|_F^2/n = T = 2n$.
With $\alpha = n/(n+T) = 1/3$ (and hence $1-\alpha = 2/3$),
equation~\eqref{eq:kappa} gives
\[
  \kappa_{\mathrm{eff}} \approx
  \frac{\tfrac{2}{3}(5.83n) + \tfrac{1}{3}(2n)}{\tfrac{2}{3}(0.17n) + \tfrac{1}{3}(2n)}
  = \frac{3.89n + 0.67n}{0.113n + 0.67n} \approx 5.8.
\]
Since~\eqref{eq:minres_conv_proxy} is a Chebyshev
proxy, exact iteration counts cannot be read off from $\kappa_{\mathrm{eff}}$ alone,
but the predicted compression is qualitatively consistent with the halving of iterations
(62 to 32) observed in Table~\ref{tab:synthetic}.  In practice, MINRES often converges
faster than the linear rate predicted by~\eqref{eq:minres_conv_proxy}: as the iteration
proceeds, eigenvalues are progressively absorbed into the polynomial, and the effective
approximation problem shrinks, leading to superlinear convergence \cite{simoncini2013}.
For real financial returns the sample covariance is far more ill-conditioned, explaining
the much higher iteration counts on S\&P~500 data (Table~\ref{tab:sp500}).
\end{remark}

\subsection{Inequality constraints and the active set}

Each active-set outer iteration replaces $(A, b)$ with the augmented pair
$A_{\mathcal{A}} = [A,\, C_{:,\mathcal{A}}]$,
$b_{\mathcal{A}} = (b^\top, d_{\mathcal{A}}^\top)^\top$ in system~\eqref{eq:kkt},
where $\mathcal{A}$ indexes the currently binding inequality constraints.
Section~\ref{sec:activeset} gives the full algorithm.

\section{Active-Set Method}\label{sec:activeset}

General linear inequalities $C^\top w \leq d$ are handled by an active-set
loop that converts the inequality-constrained problem into a sequence of
equality-constrained subproblems, each solvable by the KKT machinery of
Section~\ref{sec:kkt}.

\subsection{Algorithm}

Let $\mathcal{A} \subseteq \{1,\ldots,p\}$ denote the current active set—the
index set of inequalities treated as equalities.  At each iteration the
augmented constraint matrix $A_{\mathcal{A}} = [A,\, C_{:,\mathcal{A}}]$ and
right-hand side $b_{\mathcal{A}} = (b^\top, d_{\mathcal{A}}^\top)^\top$ replace
$A$ and $b$ in the KKT system~\eqref{eq:kkt}, and the resulting
equality-constrained subproblem is solved exactly.  The solution $w$ is then
checked against all inactive constraints: any constraint $i \notin \mathcal{A}$
with $C_{:,i}^\top w > d_i$ (up to tolerance $10^{-10}$) is promoted to
$\mathcal{A}$.  The loop terminates when no violation is found.

Promoting all violators simultaneously---rather than one at a time---reduces
the number of outer iterations in practice; correctness is preserved because the KKT
system is re-solved from scratch at each step with the full updated active set.

\begin{remark}[Termination, complexity, and infeasibility]
The active set $\mathcal{A}$ is monotonically non-decreasing: constraints are promoted
but never demoted.  For a \emph{feasible} problem, since $|\mathcal{A}| \leq p$, the loop
terminates in at most $p + 1$ outer iterations regardless of problem conditioning,
provided the inner solver is exact.  If the constraint data $(A, b, C, d)$ defines an
empty feasible set, the augmented KKT system becomes singular at some step; the direct
solver raises a \texttt{LinAlgError} and the iterative solvers may not converge.
Infeasibility is not detected explicitly: users are responsible for verifying that the
supplied constraints are consistent before calling any solver.

The worst-case total cost is $O\!\bigl(p\cdot(n+m+p)^3\bigr)$ for the direct KKT
variant, since each of the at most $p+1$ outer steps solves a system of size at most
$n + m + p$; for the long-only default $p = n$ with $m = 1$ this gives $O(n^4)$.  The
MINRES variant costs $O(p\cdot k(Tn + mn))$ in the worst case.  In practice the active
set stabilises after very few steps (e.g.\ $7$ outer iterations for $n = 495$ in
Table~\ref{tab:sp500}), so observed cost is far below the worst case.
\end{remark}

\begin{remark}[Dual feasibility and global optimality]\label{rem:dualfeas}
At termination $w$ is primal feasible.  For global optimality the KKT conditions
also require dual feasibility: the multipliers $\nu_i \geq 0$ for all $i \in
\mathcal{A}$ (note $\nu_i$ is distinct from the expected-return vector $\mu$
of~\eqref{eq:primal}).  A standard active-set QP solver enforces this by removing
constraints whose multipliers turn negative; the present loop omits that step.

Dual feasibility is nonetheless guaranteed here because every constraint is promoted
solely by violation.  To see this, fix any $i \in \mathcal{A}$ and consider the
parametric optimal value
\[
  v(\delta) = \min\bigl\{f(w) : A^\top w = b,\;
    C_{:,j}^\top w \leq d_j\;\forall j \in \mathcal{A},\;
    C_{:,j}^\top w \leq d_j\;\forall j \notin \mathcal{A}\bigr\}
\]
with the single constraint $j = i$ replaced by $C_{:,i}^\top w \leq d_i + \delta$.
Since the feasible set is non-decreasing in $\delta$, the strictly convex objective
gives $v(\delta) \leq v(0)$ for $\delta > 0$.  By the envelope theorem,
$v'(0) = -\nu_i \leq 0$, so $\nu_i \geq 0$.

This holds for every $i \in \mathcal{A}$, including constraints promoted
simultaneously: each was violated by the previous iterate, confirming the
unconstrained optimum pressed against the bound from the infeasible side, which is
precisely the condition that makes the envelope argument valid.  Combined with primal
feasibility and stationarity (which hold by construction of the KKT solve), the full
KKT conditions are satisfied at termination, confirming $w = w^*$.

\emph{Scope and benchmark verification.}
The argument requires the following condition: every constraint $i \in \mathcal{A}$
at termination was violated by the iterate at which it was first added.  A
sufficient condition is that the bound $d_i$ lies in the interior of the range of
$C_{:,i}^\top w$ over the unconstrained-budget-constraint solution set; this is
guaranteed whenever no inequality constraint is binding at the globally
unconstrained minimum.

All benchmarks in this paper satisfy this condition.  The synthetic and
efficient-frontier experiments use long-only constraints ($w_i \geq 0 \Leftrightarrow
-w_i \leq 0$) and sector caps ($\sum_{i \in S_k} w_i \leq 0.25$).  For i.i.d.\
Gaussian data with $T = 2n$ returns, the unconstrained minimum-variance portfolio
has many negative weights (by the sign symmetry of a standard Gaussian matrix),
so the long-only constraints are violated before being activated.  The sector caps
are also violated by the unconstrained portfolio with high probability.  The S\&P~500
experiments impose only the budget equality constraint ($A^\top w = \mathbf{1}^\top w = 1$)
and use no inequalities, so the active set remains empty and the dual-feasibility
argument is vacuous.  The argument does \emph{not} extend to lower bounds
$w_i \geq \ell_i > 0$ that may bind at the optimum without being violated during
the loop; for such problems a standard multiplier check and constraint-removal step
is needed.
\end{remark}

\begin{remark}[Constraint satisfaction with iterative solvers]
When an iterative solver (MINRES, CG) is used with finite tolerance, the constraint
violations used to extend $\mathcal{A}$ are computed from an approximate $w$; the
correctness argument of Remark~\ref{rem:dualfeas} still applies, but the solution
returned may not satisfy all constraints to full machine precision.  Tightening the
solver tolerance (e.g.\ to $10^{-8}$) or switching to the direct KKT solver is
recommended when exact constraint satisfaction is required.
\end{remark}

\begin{remark}
Proposition~\ref{prop:nonsingular} requires $A_{\mathcal{A}}$ to have full column rank.
This can fail if two promoted inequality columns are linearly dependent---for instance,
if one sector constraint is a strict subset of another.  In practice this is avoided by
checking that each newly promoted column is linearly independent of the current columns
of $A_{\mathcal{A}}$ before adding it; a dependent column indicates a redundant
constraint that can be dropped without altering the feasible set.
\end{remark}

\subsection{Implementation}

The inner solver is injected as a subroutine $f(\mathcal{A}) \to (w, k)$
that accepts the current active set and returns the primal solution together
with an inner iteration count, separating the active-set loop from the choice
of linear solver.  Section~\ref{sec:solvers} presents suitable implementations.

\begin{remark}\label{rem:warmstart}
Tracing the efficient frontier requires solving problem~\eqref{eq:primal} for a range
of $\rho$ values.  Because the saddle-point matrix $2((1-\alpha)X^\top X + (\alpha\|X\|_F^2/n)I)$
is independent of $\rho$, only the right-hand side changes between solves.  The
function-pointer interface would make warm-starting natural---the active set from one
$\rho$ could be used as the initial $\mathcal{A}$ for the next, potentially reducing
outer iterations---though the current experiments solve each $\rho$ from scratch;
see~\S\ref{sec:results}.
\end{remark}

\begin{algorithm}[ht]
\caption{Active-set loop}\label{alg:activeset}
\begin{algorithmic}[1]
\Require Inner solver $f$;\; inequality data $C \in \mathbb{R}^{n \times p}$, $d \in \mathbb{R}^p$;\; tolerance $\varepsilon = 10^{-10}$
\Ensure Weight vector $w$;\; total inner iteration count $k$
\State $\mathcal{A} \leftarrow \emptyset$;\; $k \leftarrow 0$
\Repeat
  \State $(w,\, k_{\mathrm{step}}) \leftarrow f(\mathcal{A})$;\; $k \leftarrow k + k_{\mathrm{step}}$
  \State $\mathcal{V} \leftarrow \{i \notin \mathcal{A} : C_{:,i}^\top w - d_i > \varepsilon\}$
  \State $\mathcal{A} \leftarrow \mathcal{A} \cup \mathcal{V}$
\Until{$\mathcal{V} = \emptyset$}
\State \Return $w$, $k$
\end{algorithmic}
\end{algorithm}

\section{Computational Methods}\label{sec:solvers}

\subsection{Direct KKT Solver}

At each active-set step the saddle-point system~\eqref{eq:kkt} is assembled
explicitly and solved by \texttt{scipy.linalg.solve} with
\texttt{assume\_a=\textquotesingle{}sym\textquotesingle{}} (Bunch--Kaufman symmetric
factorisation \cite{bunch1977}), at the cost of one dense $O\!\bigl((n+m_\mathcal{A})^3\bigr)$
factorisation per outer iteration, where $m_\mathcal{A} = m + |\mathcal{A}|$ is the
total number of active constraints.  The symmetric factorisation uses half the arithmetic of a general LU,
though at the moderate problem sizes studied here the observed wall-clock runtimes
are slower than \texttt{numpy.linalg.solve}, which dispatches to a general LU via
a more direct LAPACK path.
After convergence the weight vector is optionally clipped to $[0,\infty)$ and
renormalised to sum to one; this projection is correct for the default long-only
problem but should be disabled when custom constraint matrices are supplied.

\begin{algorithm}[ht]
\caption{Direct KKT solver}\label{alg:kkt}
\begin{algorithmic}[1]
\Require Problem data
\Ensure Weight vector $w$;\; outer active-set step count $s$
\Function{DirectSolve}{$\mathcal{A}$}
  \State Set $A_{\mathcal{A}} \leftarrow [A,\, C_{:,\mathcal{A}}]$,\; $b_{\mathcal{A}} \leftarrow (b^\top, d_{\mathcal{A}}^\top)^\top$
  \State Assemble $K_{\mathcal{A}}$ and $r_{\mathcal{A}}$ from~\eqref{eq:kkt}
  \State Solve $K_{\mathcal{A}} z = r_{\mathcal{A}}$ by dense factorisation
  \State \Return $z_{1:n}$,\; $1$  \hfill\textit{(per-step count; no inner iterations)}
\EndFunction
\State $(w,\, s) \leftarrow \textsc{ActiveSetLoop}(\textsc{DirectSolve})$
  \hfill\textit{(Algorithm~\ref{alg:activeset})}
\State \Return $w$, $s$
\end{algorithmic}
\end{algorithm}

The inner function returns $1$ per call; $s$ therefore counts active-set outer steps,
distinct from the cumulative Krylov iteration count $k$ reported by the iterative solvers.

\subsection{CVXPY Reference Implementation}

The general mean-variance problem maps directly onto a CVXPY \cite{diamond2016}
model solved by Clarabel \cite{clarabel2024}: the objective $(1-\alpha)\|Xw\|^2 + (\alpha\|X\|_F^2/n)\|w\|^2 - \rho\,\mu^\top w$
is expressed symbolically, the constraints $A^\top w = b$ and $C^\top w \leq d$
are passed unchanged, and the solver returns the primal solution together with
its interior-point iteration count.  This implementation serves as the
correctness reference against which the Krylov solvers are validated.

CVXPY \cite{diamond2016} with Clarabel \cite{clarabel2024} offers a remarkably
expressive modelling API: the same
few lines handle arbitrary convex objectives and constraint matrices without any
solver-specific tuning.  The Krylov approach trades this generality for
structure.  Four differences are worth noting.

\textit{Covariance matrix.}  When the objective is expressed as $\|Xw\|^2$ and passed
symbolically to CVXPY, CVXPY constructs the quadratic cost matrix $P = 2X^\top X$
during problem canonicalisation before dispatching to Clarabel, incurring $O(Tn^2)$
cost and $O(n^2)$ storage (Table~\ref{tab:synthetic} caption).  The Krylov solvers
avoid this entirely by working with $X$ directly via the matrix-free operator of
\S\ref{sec:matfree}.

\textit{Inequality constraints.}  Clarabel enforces $C^\top w \leq d$ via a
logarithmic barrier, keeping all iterates strictly interior to the feasible
set.  The active-set method of Section~\ref{sec:activeset} instead promotes
violated inequalities to equalities, solving a sequence of saddle-point systems
of increasing size.  Neither approach dominates: the barrier is smooth and
avoids combinatorial bookkeeping, while the active set exploits the Krylov
structure at each step and terminates in few outer iterations for
well-conditioned problems.

\textit{Per-iteration cost.}  Each Clarabel interior-point step solves a
dense $O(n^3)$ normal-equations system.  Each MINRES iteration costs $O(Tn + mn)$
(two passes through $X$ plus constraint-matrix products; $O(mn) \ll O(Tn)$ for $m \ll T$).
Each constraint-eliminated CG iteration costs $O(Tn + n^2)$ (two passes through $X$ plus the
null-space projections $Pv$ and $P^\top(\cdot)$; see \S\ref{sec:cg}).  For $T = 2n$,
MINRES and CG iterations have the same $O(n^2)$ cost, but MINRES operates on the
full $(n+m)$-dimensional system while CG on the reduced $(n - m_\mathcal{A})$-dimensional
one; the smaller system size does not imply faster convergence---the condition number
of $P^\top((1-\alpha)X^\top X + (\alpha\|X\|_F^2/n)I)P$ is comparable to $\kappa_\mathrm{eff}$---and
empirically MINRES requires fewer total iterations than CG (Table~\ref{tab:synthetic}).
For $k$ Krylov iterations the MINRES total
is $O(k(Tn + mn)) \ll O(n^3)$ when Ledoit-Wolf shrinkage keeps $k \ll n$; without
shrinkage, ill-conditioned real data can force $k \gg n$ (see~\S\ref{sec:results}).

\subsection{Matrix-Free KKT Operator}\label{sec:matfree}

The direct solver assembles $K$ explicitly, which requires forming $X^\top X$
at $O(Tn^2)$ cost and storing it at $O(n^2)$ memory.  Krylov methods only
require matrix-vector products $Kv$, which can be computed without
materialising $X^\top X$:
\[
  K \begin{pmatrix} v_w \\ v_\lambda \end{pmatrix}
  =
  \begin{pmatrix}
    2\bigl((1-\alpha)X^\top(X v_w) + \tfrac{\alpha\|X\|_F^2}{n}\, v_w\bigr) + A\, v_\lambda \\
    A^\top v_w
  \end{pmatrix}.
\]
Each application makes two passes through $X$ at $O(Tn)$ plus $O(mn)$ for the
constraint-matrix products, avoiding both the $O(Tn^2)$ formation cost and the
$O(n^2)$ storage of $X^\top X$.  For $k$ Krylov iterations the total cost is
$O(k(Tn + mn))$, compared with $O(Tn^2 + n^3)$ for the direct solve; with
Ledoit-Wolf shrinkage, $k \ll n$ in practice and the iterative solvers are
substantially faster for large $n$ (Table~\ref{tab:synthetic}); without shrinkage,
real-data ill-conditioning can increase $k$ substantially (Table~\ref{tab:sp500}).

The recurrence above is implemented as a matrix-free operator that MINRES and
CG accept directly; no explicit matrix is ever formed or stored.

\subsubsection{MINRES}\label{sec:minres}

MINRES \cite{paige1975} solves symmetric indefinite systems $Kx = b$ via
short recurrences, making it well-suited to the saddle-point
structure~\eqref{eq:kkt}.  At each active-set step the matrix-free operator
from \S\ref{sec:matfree} is passed directly to MINRES; the total iteration count accumulates
across all steps.  Ledoit-Wolf shrinkage ($\alpha > 0$) compresses the spectrum of $X^\top X$ and substantially
reduces the iteration count.  The current implementation uses no preconditioner;
well-studied block-diagonal preconditioners for symmetric saddle-point systems
\cite{benzi2005} could reduce iteration counts further, particularly for
ill-conditioned real data, and are left for future work.

\begin{algorithm}[ht]
\caption{MINRES solver}\label{alg:minres}
\begin{algorithmic}[1]
\Require Problem data
\Ensure Weight vector $w$;\; total MINRES iteration count $k$
\Function{MinresSolve}{$\mathcal{A}$}
  \State Build matrix-free operator $\mathcal{K}_{\mathcal{A}}$ and right-hand side $r_{\mathcal{A}}$
  \State Solve $\mathcal{K}_{\mathcal{A}} z = r_{\mathcal{A}}$ via MINRES \cite{paige1975}
  \State \Return $z_{1:n}$,\; iteration count
\EndFunction
\State $(w,\, k) \leftarrow \textsc{ActiveSetLoop}(\textsc{MinresSolve})$
  \hfill\textit{(Algorithm~\ref{alg:activeset})}
\State \Return $w$, $k$
\end{algorithmic}
\end{algorithm}

\subsubsection{Conjugate Gradient in the Null Space}\label{sec:cg}

An alternative is to eliminate the equality constraints via QR factorisation of
$A_{\mathcal{A}}$, project onto $\ker(A_{\mathcal{A}}^\top)$, and apply CG
\cite{hestenes1952} to the resulting $(n - m_{\mathcal{A}})$-dimensional SPD system.
The per-iteration cost is $O(Tn + n^2)$ versus $O(Tn + mn)$ for MINRES, and the
one-time QR setup costs $O(n^2 m_{\mathcal{A}})$ per active-set step to form the
full orthogonal factor---reaching $O(n^3)$ when $m_{\mathcal{A}} = O(n)$.  For
general $A$ this QR cost dominates and makes the constraint-eliminated CG slower
than MINRES; the benchmarks in \S\ref{sec:results} confirm this.

\section{Numerical Results}\label{sec:results}

All experiments were run on an Apple M4 Pro (14-core CPU, 48\,GB RAM) under
Python~3.12, NumPy~2.4, SciPy~1.17, CVXPY~1.8.2, and Clarabel~0.11.1.
Iterative solvers use SciPy's default relative residual tolerance of $10^{-5}$;
reported times are the minimum of three repeated solves.  Speedup figures are
hardware- and version-dependent; the dominant factor is the $O(Tn^2)$
problem-canonicalisation cost incurred by CVXPY~1.8.2 when constructing the
quadratic cost matrix $P = 2X^\top X$ before dispatching to Clarabel---future
CVXPY releases may reduce this overhead.  All implementations are available in
the companion package \cite{fast_minimum_variance}.

\subsection{Synthetic benchmark}\label{sec:synth}

Table~\ref{tab:synthetic} reports results for a problem with $n=1000$ assets,
$T=2000$ i.i.d.\ Gaussian returns (generated with \texttt{numpy.random.default\_rng(0)}),
five equal-sized sector caps at $25\%$, return-tilt weight $\rho=0.5$, and
expected-return vector $\mu\sim\mathcal{N}(0,I)$.  Two panels show results without
and with Ledoit-Wolf shrinkage; the LW panel uses $\alpha = n/(n+T) = 1/3$.

\begin{table}[ht]
\centering
\begin{tabular}{lccc}
\toprule
Method & Time (s) & Iterations & Speedup \\
\midrule
cvxpy (Clarabel)             & 10.77 & n/a &    1.0x \\
KKT direct                   & 0.087 &  --$^*$ &  124.1x \\
MINRES                        & 0.010 &  62 & 1046.6x \\
CG (constr.-elim.)            & 0.104 & 105 &  103.7x \\
\midrule
MINRES + LW                   & 0.008 &  32 & 1325.8x \\
CG + LW                       & 0.029 &  39 &  366.7x \\
\bottomrule
\end{tabular}
\caption{Synthetic benchmark: $n=1000$ assets, $T=2000$ i.i.d.\ Gaussian returns,
five sector caps at $25\%$, $\rho=0.5$, $\mu\sim\mathcal{N}(0,I)$.  Lower panel adds
Ledoit-Wolf shrinkage ($\alpha=1/3$).  Speedup for the upper panel is relative to
\texttt{cvxpy} ($10.77$~s); speedup for the lower panel is also relative to the same
\texttt{cvxpy} baseline, as applying LW to \texttt{cvxpy} does not measurably change
its runtime (the ridge coefficient $\alpha\|X\|_F^2/n$ is passed as a scalar to
Clarabel and the interior-point complexity remains $O(n^3)$ per step).  The reported \texttt{cvxpy} times include the $O(Tn^2)$
cost of forming $X^\top X$: the benchmark passes $X$ directly and CVXPY constructs
the quadratic cost matrix $P = 2X^\top X$ during problem canonicalisation before
dispatching to Clarabel.
The Iterations column counts cumulative \emph{inner} Krylov iterations across all
active-set steps for MINRES and CG; \texttt{cvxpy} interior-point iterations are not
comparable (each costs $O(n^3)$ versus $O(Tn + mn)$ per Krylov step) and are therefore
omitted.  $^*$KKT uses a direct factorisation at each active-set step and has no inner
iterations; the $4$ \emph{outer} active-set steps taken are distinct from inner
iterations and are not shown in this column.  The small number of outer steps
reflects the batch-promotion strategy: all violated constraints (long-only and
sector caps) are added to the active set simultaneously at each pass, so
feasibility is typically achieved in fewer than five iterations.}
\label{tab:synthetic}
\end{table}

Without shrinkage, MINRES is the fastest solver at $0.010$~s ($1047\times$ faster
than CVXPY), outperforming KKT direct by ${\approx}9\times$.  The constraint-eliminated CG
formulation, however, is \emph{slower} than the direct solve: $0.104$~s versus
$0.087$~s.  This is a qualitative difference from the minimum variance setting of
\cite{minvar_paper}, where CG was competitive with MINRES.  The cause is the $O(n^2 m_\mathcal{A})$ cost of forming the full orthogonal factor $Q$
via QR factorisation of the active constraint matrix (see \S\ref{sec:cg}), which
reaches $O(n^3)$ when $m_\mathcal{A} = O(n)$.  For the budget-only constraint
$A = \mathbf{1}$ the companion paper avoids this by exploiting an implicit $O(n)$
Householder reflector; for general $A$ no such shortcut is available and the QR
dominates at large $n$.

Ledoit-Wolf shrinkage compresses the spectrum and reduces MINRES from 62 to 32
iterations, cutting runtime to $0.008$~s ($1326\times$ speedup).  CG also benefits
from shrinkage (105 to 39 iterations), but the QR bottleneck limits the gain:
CG+LW reaches $0.029$~s, still ${\approx}4\times$ slower than MINRES+LW.

\subsection{Scaling with problem size}\label{sec:scaling}

Figure~\ref{fig:scaling} shows wall-clock time as a function of $n$ for the
constrained Markowitz problem (five sector caps, $\rho = 0.5$, LW shrinkage, $T=2n$).
The KKT direct solve scales as $O(n^3)$, dominated by the dense $n \times n$ linear
system.  MINRES+LW scales close to $O(n^2)$: each iteration costs $O(Tn) = O(n^2)$
for $T = 2n$, and the Krylov iteration count $k$ is approximately constant in $n$ under LW
shrinkage---since $\kappa_{\mathrm{eff}} \approx 5.8$ for $T = 2n$ Gaussian data regardless
of $n$ (Remark~\ref{rem:marchenko}), the convergence bound~\eqref{eq:minres_conv_proxy} predicts a fixed
iteration count---so the total cost $O(kTn)$ follows slope 2 in the log--log plot of
Figure~\ref{fig:scaling}.
MINRES+LW opens a visible gap above $n \approx 200$.  CG+LW is faster than KKT for small $n$ but its scaling is degraded by the
$O(n^2 m_\mathcal{A})$ QR cost for general $A$.

\begin{figure}[ht]
\centering
\includegraphics[width=\linewidth]{graphs/markowitz_scaling.pdf}
\caption{Left: wall-clock time vs.\ $n$ (log--log, $T=2n$, five sector caps,
$\rho=0.5$, LW shrinkage throughout).  Right: efficient frontier for $n=500$ assets
with five sector caps, $\alpha=1/3$, tracing 51 risk--return points as $\rho$ varies
from $0$ to $2$, computed with the KKT direct solver ($\alpha=1/3$) for illustration;
MINRES+LW and CG+LW (same $\alpha$) produce the same curve to relative weight
tolerance $10^{-4}$ over all 51 points (see \S\ref{sec:frontier}).}
\label{fig:scaling}
\end{figure}

\subsection{Efficient frontier}\label{sec:frontier}

Tracing the full efficient frontier requires solving the Markowitz problem for a
range of $\rho$ values.  For $n=500$ assets, $T=1000$ returns ($T = 2n$, matching
the scaling experiments), and five sector caps, we time 21 equally spaced points at
$\rho \in [0, 2]$ under two scenarios: (i)~no shrinkage ($\alpha=0$, the same
optimisation problem for every solver), and (ii)~Ledoit-Wolf shrinkage
($\alpha = n/(n+T) = 1/3$, applied identically to every solver including CVXPY).

\textit{Algorithmic comparison ($\alpha=0$, same problem).}
KKT direct takes $0.393$~s ($18.7$~ms per point), a $91\times$ speedup over
CVXPY at $36.0$~s ($1713$~ms); MINRES without shrinkage takes $0.153$~s
($7.3$~ms, $235\times$).

\textit{Regularised comparison ($\alpha=1/3$, same problem).}
Applying LW shrinkage to CVXPY does not measurably change its runtime:
$36.0$~s~$\to$~$36.1$~s (interior-point cost is dominated by factorisation, not
problem structure).  MINRES+LW takes $0.043$~s ($2.0$~ms per point,
${\approx}838\times$ speedup over CVXPY+LW), and CG+LW takes $0.155$~s
($7.4$~ms, ${\approx}233\times$).  MINRES+LW is the fastest by
${\approx}3.6\times$ over CG+LW.

The Krylov iteration count per frontier point is approximately constant in $n$
(Remark~\ref{rem:marchenko}), consistent with the single-$\rho$ figures of
Table~\ref{tab:synthetic} at $n=1000$, confirming that repeated solves across
the frontier do not inflate the Krylov work.  The right panel of
Figure~\ref{fig:scaling} shows the risk--return curve computed by the KKT direct
solver ($\alpha=0$) using 51 equally spaced $\rho$ values (denser than the
21-point timing experiment, for a smoother illustration).  MINRES+LW and CG+LW ($\alpha=1/3$)
produce the same 51-point frontier as KKT direct ($\alpha=1/3$), verified by
comparing portfolio weights in the $\ell_\infty$ norm to relative tolerance $10^{-4}$
at every frontier point.  The saddle-point matrix
$2\bigl((1-\alpha)X^\top X + \tfrac{\alpha\|X\|_F^2}{n}I\bigr)$ is independent of
$\rho$, so the active-set and Krylov work is amortised across frontier points
(see Remark~\ref{rem:warmstart} for the warm-starting implication).

\subsection{S\&P~500 real-data benchmark}

Table~\ref{tab:sp500} reports results on $n=495$ S\&P~500 constituents over $T=1192$
daily log-returns (July 2021--April 2026) sourced from Yahoo Finance via the
\texttt{yfinance} Python package; the data construction script will be made available
in the companion package \cite{fast_minimum_variance} upon publication.  The return matrix $X$ is formed by
demeaning the log-returns column-wise (subtracting each asset's sample mean over the
estimation window).  The LW panel uses $\alpha = n/(n+T) = 495/1687 \approx 0.293$.
This setting is the base case of \cite{minvar_paper}
and is chosen deliberately to isolate the effect of real equity cross-correlations and
heavy tails on solver performance, independent of the sector-cap constraint structure
tested in \S\ref{sec:synth}--\S\ref{sec:scaling}.  A natural extension---applying sector
caps and a return tilt to real data---is left for future work.  The sample covariance is
significantly more ill-conditioned than the Gaussian benchmark.  The column $\|Xw\|$
in Table~\ref{tab:sp500} is the square root of the sample portfolio variance scaled by
$T$: $\|Xw\|^2 = T\,\hat\sigma_p^2$ where $\hat\sigma_p^2 = w^\top\hat\Sigma w$ and
$\hat\Sigma = X^\top X / T$ is the sample covariance.  Within each panel (fixed
shrinkage setting) a smaller value indicates lower in-sample portfolio variance;
the two panels optimise different objectives and their $\|Xw\|$ values are not
directly comparable across panels.

\begin{table}[ht]
\centering
\begin{tabular}{lcccc}
\toprule
Method & $\|Xw\|$ & Time (s) & Iterations & Speedup \\
\midrule
\multicolumn{5}{l}{\textit{Without Ledoit-Wolf shrinkage}} \\[2pt]
cvxpy (Clarabel)                    & 0.2241 & 1.614 &  n/a &  1.0x \\
KKT direct                          & 0.2246 & 0.076 & --$^*$ & 21.2x \\
MINRES                              & 0.2255 & 0.050 & 1065 & 32.3x \\
CG (constr.-elim.)                  & 0.2246 & 0.081 &  385 & 19.9x \\
\midrule
\multicolumn{5}{l}{\textit{With Ledoit-Wolf shrinkage}} \\[2pt]
cvxpy (Clarabel)                    & 0.2281 & 1.488 &  n/a &  1.0x \\
KKT direct                          & 0.2280 & 0.080 & --$^*$ & 18.5x \\
MINRES + LW                         & 0.2286 & 0.021 &  424 & 70.8x \\
CG + LW (constr.-elim.)             & 0.2280 & 0.069 &  117 & 21.6x \\
\bottomrule
\end{tabular}
\caption{S\&P~500 universe: $n=495$ assets, $T=1192$ trading days
(July~2021--April~2026), long-only minimum-variance problem.
Speedup relative to \texttt{cvxpy} within each panel.
Iterations are cumulative Krylov iterations across all active-set steps;
\texttt{cvxpy} interior-point iterations are not comparable to Krylov iterations
(different cost per step) and are therefore omitted.
$^*$KKT is a direct inner solve; the 7 active-set outer steps (each promoting a batch
of assets whose unconstrained weight is negative to the long-only bound $w_i = 0$)
are not shown.}
\label{tab:sp500}
\end{table}

As a correctness check, the $\|Xw\|$ values in Table~\ref{tab:sp500} confirm
that all four solvers agree on in-sample portfolio risk (measured in $\|Xw\|$,
the square root of the sample portfolio variance scaled by $T$) to within
$(0.2255 - 0.2241)/0.2241 \approx 0.6\%$ relative error in $\|Xw\|$, or
equivalently $\approx 1.2\%$ in portfolio variance $\|Xw\|^2$, relative to
the CVXPY reference.  The deviation is entirely attributable to MINRES: KKT
direct and CG agree with CVXPY to within $0.2\%$ in $\|Xw\|$.  At the stated
solver tolerance of $10^{-5}$ on the relative residual, a $0.6\%$ deviation in
$\|Xw\|$ is expected for high-condition-number data: for practical minimum-variance
applications the dominant source of out-of-sample error is estimation noise in $X$,
which dwarfs the $0.6\%$ in-sample solver discrepancy; tightening to $10^{-7}$
would reduce it to negligible levels at the cost of a few hundred additional
iterations (see the convergence discussion below).

Without shrinkage, MINRES requires 1065 iterations---$17\times$ more than on the
Gaussian benchmark---reflecting the much higher condition number~\eqref{eq:kappa} of
real return data; despite the large iteration count, MINRES is nonetheless faster than
KKT direct in this regime ($32\times$ versus $21\times$ speedup), because the Bunch--Kaufman
factorisation overhead of the direct solve is not amortised at this problem size.
The MINRES portfolio risk ($\|Xw\| = 0.2255$) is slightly higher than for
the other solvers ($0.2246$): iterative termination at a relative residual of $10^{-5}$
leaves a small inexactness in the KKT conditions, so the returned weight vector is not
the exact minimiser.  From the convergence bound~\eqref{eq:minres_conv_proxy}, each additional
decade of accuracy requires $O\!\bigl(1/\log\tfrac{\sqrt{\kappa_{\mathrm{eff}}}+1}{\sqrt{\kappa_{\mathrm{eff}}}-1}\bigr)$
further iterations; for the high-$\kappa$ S\&P~500 data this is several hundred additional
steps per decade, making tight tolerances expensive without shrinkage.  Ledoit-Wolf shrinkage reduces the MINRES count to 424 and lifts
its speedup to $71\times$, making it the fastest solver.  The slightly higher portfolio
risk under LW ($\|Xw\| = 0.2280$--$0.2286$ versus $0.2241$--$0.2246$ without LW) is
expected: LW adds a ridge term $\tfrac{\alpha\|X\|_F^2}{n}\|w\|^2$ to the objective that penalises
extreme weights, so the resulting portfolio sacrifices a small amount of in-sample variance
reduction in exchange for a more stable estimator.  Even after shrinkage, the
iteration count is $13\times$ larger than on Gaussian data, confirming that the low
iteration counts in \S\ref{sec:synth} reflect the near-identity covariance of
i.i.d.\ Gaussian returns rather than a general property of the method.

\section{Conclusion}

The Markowitz mean-variance problem, formulated directly in terms of the returns
matrix, inherits the KKT saddle-point structure identified in \cite{minvar_paper}
and admits the same MINRES treatment without modification.  Four generalisations
of the minimum-variance setting are required---general equality constraints,
general linear inequalities via active-set, a return tilt, and Ledoit-Wolf
shrinkage---and each preserves the symmetric indefinite structure that MINRES
exploits.  No explicit covariance matrix is ever formed.

The paper also describes a modular software architecture in which the
active-set loop is decoupled from the inner linear solver via a function-pointer
interface (\S\ref{sec:activeset}), and the KKT saddle-point operator is exposed as a
\texttt{LinearOperator} (\S\ref{sec:matfree}) that MINRES and CG consume without modification.
The full problem specification is encapsulated in a single frozen \texttt{Problem}
dataclass (\S\ref{sec:problem}), providing a clean API for all solvers.

The empirical results confirm that the speedup from \cite{minvar_paper} extends
to the full constrained Markowitz setting.  MINRES with Ledoit-Wolf shrinkage is
$1326\times$ faster than CVXPY on the synthetic benchmark ($n=1000$) and traces
the efficient frontier for $n=500$ assets in $0.043$~s versus $36.1$~s for
CVXPY with the same shrinkage (${\approx}838\times$ speedup; both solvers use
$\alpha=1/3$).  The
constraint-eliminated CG formulation \cite{hestenes1952} incurs an $O(n^2 m_\mathcal{A})$
cost from the QR null-space basis for general $A$ (reaching $O(n^3)$ when many
constraints are active), making it slower than MINRES for large $n$; restoring the
$O(n)$ implicit Householder efficiency for the general case is a direction for future
work, as is warm-starting the active set across efficient frontier points.  The dual-feasibility guarantee (Remark~\ref{rem:dualfeas}) is verified to hold for
all reported benchmarks; the remark characterises the class of problems for which
the add-only strategy is correct and identifies the extension to lower-bounded
weights as the remaining scope limitation.

\section*{Acknowledgements}

The authors thank Martin~H.~Gutknecht (ETH~Z\"{u}rich) for supervising the 2004 diploma
thesis \cite{schmelzer2004} on block MINRES, which motivated the present work.

\bibliographystyle{siam/siamplain}
\bibliography{bib/refs}

\end{document}
